\section*{Acknowledgements}

We thank S. Petrarca for an early participation to this work and for
many discussions. We also thank E. Gabrielli and the members
of LPTHE
for discussions. G.M. and C.P. thank the theory division at CERN,
where part of this study has been performed. We acknowledge
the
partial support by  M.U.R.S.T. and by the EC contract CHRX-CT92-0051.
C.P. acknowledges the support by the Human Capital and Mobility
Program, contract  ERBCHBICT930887.
CTS acknowledges the Particle Physics and Astronomy Research Council
for its support through the award of a Senior Fellowship.

\begin{figure}[htbp]
\centering
\includegraphics[width=0.62\textwidth]{images/file.pdf}
\caption{Diagrams with standard, a),  and scalar quarks, b).}
\label{fig:appe}
\end{figure}

\section*{Appendix}

In this appendix we demonstrate that the non-perturbative contributions
to forward two-quark matrix elements
are suppressed as inverse powers of $p^2$, where $p$ is the external momentum.
Let us consider the non-amputated, two-fermion Green function
\beq F_O(p)=\int d^4x d^4y e^{-i p \cdot x} \langle \bar q(0) \Gamma
q(x) O_\Gamma(y) \rangle ,\nn \eeq
where $O_\Gamma(y)=\bar q(y) \Gamma q(y)$ and $\Gamma$
is one of the Dirac matrices. It is natural to worry about the infrared
behaviour of  $F_O(p)$, because it contains the insertion
of a zero-momentum operator.  On general
grounds,  we expect infrared divergences to appear in Green functions computed
at exceptional momenta.  In  physical terms, this is due to the fact that
the matrix elements of $O_\Gamma$, at zero momentum, carry
information about the physical mass spectrum, which
is unaccessible to perturbation theory. In  particular, if
the  operator $O_\Gamma$  is the pseudoscalar
density ($\gamma_5$) and we work in
the chiral limit,  a contribution from  a Goldstone pole, corresponding to the
pion, should appear in $F_O(p)$. The presence of non-perturbative contributions
to $F_O(p)$ is signaled by the appearence of infrared divergences
in the  perturbative expansion. This is what happens in general. For quark
degrees of freedom in the limit of large $p^2$ however, the Operator Product
Expansion  (OPE) guarantees that $F_O(p)$ can be reliably  computed  at all
orders in perturbation theory. The argument goes as follows. In the limit of
large $p^2$, the dominant contribution to $F_O(p)$
comes from  regions of integration (in $x$ and $y$) where
the integrand is singular. This implies that, in this limit, we have to study
the behaviour of
\beq F_O(x,y)=\langle \bar q(0) \Gamma q(x) O_\Gamma(y)  \rangle,\nn \eeq
when $x \sim 0$ and $x \sim y$. In both
cases  of course, we can use the OPE in the following form\footnote{From a
technical point of view, OPE is a weak operator statement and, when $O_\Gamma$
is inserted in a Green function containing other operators, a slightly modified
form would be appropriate. It is easy to convince oneself, however, that the
additionals terms, required in the general case, are not relevant to the
present
discussion.}
\beq q_\alpha (x) O_\Gamma (y) \to A^\Gamma_{\alpha\beta}(x-y)
 \Gamma_{\beta\delta} q_\delta (y)+\dots, \nn \eeq
for $x \sim y$, and
\beq \bar q(x) \Gamma q(0) \to B_\Gamma (x^2) \bar q(0) \Gamma q(0) +\dots
,\nn \eeq
for $x \sim 0$.

Since QCD is  asymptotically free, the
 behaviour of $A^\Gamma_{\alpha\beta}(x-y) $ and $B_\Gamma (x)$ can be
computed
in perturbation theory. The interaction gives rise to logarithmic
corrections to the free field behaviour
\beq  A_\Gamma (x) =\xslash C_\Gamma (x^2) \equiv \xslash \frac{a_\Gamma
\ln^{\alpha_\Gamma}(\mu^2 x^2)+\dots}{x^2} \nn \eeq
\beq B_\Gamma(x^2)= b_\Gamma \ln^{\beta_\Gamma}(\mu^2 x^2)+\dots \nn \eeq
Using the above equations, in the limit $p^2 \to \infty$, we find
\beq \int d^4x d^4y e^{-i p \cdot x} \left( \langle \bar q(0)\Gamma  A_\Gamma
 (x-y) \Gamma
q(y) \rangle + B_\Gamma(x^2) \langle \bar q(0) \Gamma q(0) O_\Gamma(y) \rangle
 \right) = \nn \eeq \beq -i p_\mu \bar C_\Gamma (p^2) \bar S^\mu_\Gamma (p)+
\bar B_\Gamma (p^2) \bar \Delta_\Gamma (0) ,\nn \eeq
where \beq \bar S^\mu_\Gamma (p)= \int d^4y e^{-i p \cdot y}\langle \bar q(0)
 \Gamma
\gamma^\mu \Gamma q(y) \rangle
\,\,\,\,\,\,\,\,\,\,\,\,\,
 \bar \Delta_\Gamma (0)= \int d^4y \langle \bar q(0) \Gamma q(0)
O_\Gamma(y) \rangle . \nn \eeq

The term $-i p_\mu \bar C_\Gamma (p^2) \bar S^\mu_\Gamma (p)$ is
computable in perturbation theory. On
 dimensional ground, for $p^2 \to \infty$, it behaves as
\beq -i p_\mu \bar C_\Gamma (p^2) \bar S^\mu_\Gamma (p) \to
c_\Gamma \frac{\ln^{\gamma_\Gamma}(p^2/\mu^2)}{p^2} . \nn \eeq
 As for $\bar B_\Gamma (p^2)$, we find
\beq   \bar B_\Gamma (p^2) \to
d_\Gamma \frac{\ln^{\delta_\Gamma}(p^2/\mu^2)}{p^4} . \nn \eeq
$\bar B_\Gamma (p^2)$   is, however, multiplied by $\bar \Delta_\Gamma (0)$,
 which is a  non-perturbative quantity. We have shown  however,
that, in the large   $p^2$  limit, the perturbative contribution dominates by
one  power of $p^2$ over the non-perturbative one. This  is the
 reason why, even at an exceptional external  momentum,
 $F_O(p)$ is infrared safe  in perturbation theory, when $p^2$ is
 large. It is straightforward to trace the suppression
of the infrared divergence, by analyzing the corresponing Feynman
diagrams.
For example, the infrared divergent part of the
diagram in  fig.\ \ref{fig:appe}-a is suppressed by one power of $p^2$,
due to the momentum flow through the lower gluon line.

It is interesting to see what would happen in a
(fictitious)
theory, with quarks  represented  by scalar fields. In this case
the infrared divergence is not supressed, as can be seen from
the diagram in fig.\ \ref{fig:appe}-b.  The general argument goes as follows
\beq  F(p^2) =\int d^4x d^4y e^{-i p \cdot x} \langle \bar \phi (0) \phi(x)
J(y) \rangle, \nn \eeq  with \beq J(y)= \bar \phi (y) \phi(y). \nn \eeq
 When $x \sim 0$ and $x \sim y$, the OPE  gives
\beq  \phi (x) J(y) \to A\left((x-y)^2\right) \phi (y) + \dots
\,\,\,\,\,\,\,\,\,\,\,\,\,
 \bar \phi (x) \phi (0) \to B(x^2) \bar \phi (0) \phi (0) + \dots, \nn \eeq
 with
\beq  A (x^2) =  \frac{a
\ln^{\alpha}(\mu^2 x^2)+\dots}{x^2}
\,\,\,\,\,\,\,\,\,\,\,\,\,
B(x^2)= b \ln^{\beta}(\mu^2 x^2)+\dots, \nn \eeq
and we would get
\beq F(p^2) \to \bar A(p^2) \bar S(p^2) + \bar B(p^2) \bar \Delta(0) \nn \eeq
with
\beq \bar S (p^2)= \int d^4y e^{-i p \cdot y}\langle \bar \phi(0)
 \phi(y) \rangle
\,\,\,\,\,\,\,\,\,\,\,\,\,
 \bar \Delta (0)= \int d^4y \langle \bar \phi(0)  \phi(0)
J(y) \rangle . \nn \eeq
One now obtains
\beq  \bar A(p^2)  \bar S (p^2) \to
c \frac{\ln^{\gamma}(p^2/\mu^2)}{p^4}  \nn \eeq
and
\beq   \bar B (p^2) \to
d \frac{\ln^{\delta}(p^2/\mu^2)}{p^4} . \nn \eeq
The latter equations imply that non-perturbative and perturbative contributions
have the same $p^2$ dependence. Thus, we expect the appearance of infrared
divergences in $F(p^2)$.

\begin{thebibliography}{99}
\bibitem{bo} M. Bochicchio, L. Maiani, G. Martinelli, G.C. Rossi
and M. Testa, Nucl. Phys. B262 (1985) 331.
\bibitem{ma} L. Maiani and G. Martinelli, Phys. Lett. B178 (1986) 265.
\bibitem{wi} G. Martinelli, S. Petrarca, C.T. Sachrajda
and A. Vladikas, Phys. Lett. B311 (1993) 241.
Rather than quoting the numerical results of ref.[3], obtained
on a sample of 18 gauge configurations, we have opted for those
of the updated analysis of ref.\ \cite{gribov}, done with 36 gauge
configurations.
\bibitem{dallas} G. Martinelli et al.,
presented by C.T. Sachrajda, Nucl. Phys.B (Proc. Suppl.)
34 (1994) 507.
\bibitem{ber1} C. Bernard et al.,
Nucl. Phys. B (Proc. Suppl.) 17 (1990) 593; Nucl. Phys. B (Proc. Suppl.)
20 (1991) 410.
\bibitem{te} L. Maiani et al.,
Phys. Lett. B176 (1986) 445; Nucl. Phys. B289 (1987) 505
\bibitem{gav1} M.B. Gavela et al.,
Phys. Lett. B211 (1988) 139.
\bibitem{mms} L. Maiani, G. Martinelli and C.T. Sachrajda,
Nucl. Phys. B368 (1992) 281.
\bibitem{martsukuba} G. Martinelli,
 Nucl. Phys. B (Proc. Suppl.) 26 (1992) 31.
\bibitem{lepage}
G.P. Lepage and P.B. Mackenzie, Nucl. Phys. B (Proc. Suppl.) 20
(1991) 173; Phys. Rev. D48 (1993) 2250.
\bibitem{effi} M.A. Shifman and M.B. Voloshin,
Sov. J. Nucl. Phys. 45 (1987) 292; 47 (1988) 511.
\bibitem{powi} H.D. Politzer and M.B. Wise,
\pl{B206} (1988) 681; \pl{208} (1988) 504.
\bibitem{eh} E. Eichten and B. Hill, \pl {B240} (1990) 193.
\bibitem{bg} B. Grinstein, Nucl. Phys. B 339 (1990) 253.
\bibitem{efff} H. Georgi, Phys. Lett. B240 (1990) 447.
\bibitem{slami} B. Sheikholeslami and R. Wohlert, Nucl. Phys.
B259 (1985) 572.
\bibitem{altarelli} G. Altarelli, G. Curci, G.Martinelli
and S. Petrarca, Nucl. Phys. B187 (1981) 461.
\bibitem{buras} A.J. Buras, M. Jamin, M.E. Lautenbacher and P.H. Weisz,
Nucl. Phys.  B370 (1992) 69, Addendum ibid. Nucl. Phys.  B375
(1992) 501;
\noindent
A.J. Buras, M. Jamin and M.E. Lautenbacher, Nucl. Phys.  B400 (1993) 37 and
B40.
\bibitem{ciu} M. Ciuchini, E. Franco, G. Martinelli and L. Reina,
Phys. Lett. B301 (1993) 263; Nucl. Phys.  B415 (1994) 403.
\bibitem{altri} C. Parrinello, S. Petrarca and
A. Vladikas, Phys. Lett. B268 (1991) 236;
M. Paciello et al.,
Phys. Lett. B276 (1992) 163; Phys. Lett. B281 (1992) 417;
Phys. Lett. B289 (1992) 405.
\bibitem{gribov} M. Paciello et al., Rome prep. 1034 (July 1994),
to appear in Phys. Lett. B.
\bibitem{ka} L.H. Karsten and J. Smit, Nucl. Phys. B183
(1981) 103.
\bibitem{4ferm}   G. Martinelli, Phys. Lett. B141 (1984) 395; \\
C. Bernard, T. Draper and A. Soni,
 Phys. Rev. D36 (1987) 3224; \\
R. Frezzotti, E. Gabrielli, C. Pittori and
G.C. Rossi, Nucl. Phys. B373 (1991) 781;
Nucl. Phys. B409 (1993) 382.
\bibitem{msv} G. Martinelli, C.T. Sachrajda
and A. Vladikas, Nucl. Phys. B358 (1991) 212.
\bibitem{noi} G. Heatlie, G. Martinelli, C. Pittori,
G.C. Rossi and C.T. Sachrajda, Nucl. Phys. B352 (1991) 266.
\bibitem{tass2} G. Martinelli, C.T. Sachrajda, G. Salina
and A. Vladikas, Nucl. Phys. B378 (1992) 591.
\bibitem{newi} A. Borrelli, R. Frezzotti, E. Gabrielli and C. Pittori,
Nucl. Phys. B409 (1993) 382
\bibitem{hv} 't Hooft and M. Veltman, Nucl. Phys. B44 (1972) 189.
\bibitem{2ferm}
G. Martinelli and Y.C. Zhang, Phys. Lett. B123 (1983) 433; \\
E. Gabrielli, G. Heatlie, G. Martinelli, C. Pittori and
C.T. Sachrajda, Nucl. Phys. B362 (1991) 475.
\bibitem{giusti} L. Giusti and M. Testa, in preparation.
\bibitem{curci} G. Curci, Phys. Lett. B167 (1986) 425.

\end{thebibliography}
